% BHU PYQ -- Manually transcribed & error-corrected
% Paper: MTB-605 : Vector and Tensor Analysis
% Exam: B.A./B.Sc. (Hons) Semester VI Examination, 2022-23

\documentclass[11pt,a4paper]{article}
\usepackage[a4paper,top=2.5cm,bottom=2.5cm,left=2.8cm,right=2.8cm,headheight=14pt]{geometry}
\usepackage{amsmath,amssymb}
\usepackage[T1]{fontenc}
\usepackage[utf8]{inputenc}
\usepackage{lmodern}
\usepackage{microtype}
\usepackage{fancyhdr}
\usepackage{enumitem}
\usepackage{hyperref}
\usepackage{lastpage}
\usepackage{parskip}

\hypersetup{colorlinks=true,urlcolor=black,linkcolor=black,
  pdftitle={MTB-605: Vector and Tensor Analysis -- BHU PYQ},pdfauthor={Banaras Hindu University}}

\pagestyle{fancy}\fancyhf{}
\renewcommand{\headrulewidth}{0.4pt}\renewcommand{\footrulewidth}{0.4pt}
\fancyhead[L]{\small   \textit{Banaras Hindu University}}
\fancyhead[R]{\small   \textit{MTB-605: Vector and Tensor Analysis}}
\fancyfoot[C]{\small Page  \thepage\ of \pageref{LastPage}}


\newlist{parts}{enumerate}{2}
\setlist[parts,1]{label=(\alph*),leftmargin=*,itemsep=5pt,topsep=3pt}
\setlist[parts,2]{label=(\roman*),leftmargin=*,itemsep=3pt,topsep=2pt}

\newcommand{\pts}[1]{\hfill   \textit{[#1]}}


\begin{document}


\begin{center}
  {\large\bfseries BANARAS HINDU UNIVERSITY}\\[2pt]
  {\normalsize B.A./B.Sc. (Hons.) --- Semester VI Examination, 2022-23}\\[6pt]
  \rule{0.8\linewidth}{0.5pt}\\[6pt]
  {\large\bfseries Paper No. MTB-605: Vector and Tensor Analysis}\\[4pt]
  {\normalsize Time: 3 Hours \hfill Maximum Marks: 70}
\end{center}
\rule{\linewidth}{0.4pt}
\smallskip



\noindent   \textit{Instructions: Answer   \textbf{five} questions including Question~1 which is compulsory. Symbols have their usual meanings.}
\bigskip




\noindent   \textbf{Question 1.}   \textit{(Compulsory --- 2 marks each.)}

\begin{parts}
  \item If $\phi(x,y,z) = 2yz^2$ and $\vec{A} = xy\vec{i} + y^2 z\vec{j} + xz\vec{k}$, obtain $\operatorname{div}(\phi \vec{A})$. \pts{2}
  \item Given $\vec{r} = x\vec{i} + y\vec{j} + z\vec{k}$, show that $
abla(r^n) = n r^{n-2}\vec{r}$. \pts{2}
  \item Give the definition of divergence $
abla \cdot \vec{V}$ of a smooth vector field $\vec{V}(x,y,z) = V_1\vec{i} + V_2\vec{j} + V_3\vec{k}$. Obtain $
abla \cdot \vec{V}$ for $\vec{V} = yz\vec{i} + zx\vec{j} + xy\vec{k}$. \pts{2}
  \item Give the definition of curl $
abla  \times \vec{V}$ of a smooth vector field $\vec{V}(x,y,z) = V_1\vec{i} + V_2\vec{j} + V_3\vec{k}$. Obtain $
abla  \times \vec{V}$ for $\vec{V} = x^2\vec{i} + y^2\vec{j} + z^2\vec{k}$. \pts{2}
  \item Given $\vec{r} = x\vec{i} + y\vec{j} + z\vec{k}$, prove that $
abla^2(1/r) = 0$ for $r 
eq 0$. \pts{2}
  \item Give the definitions of conservative and irrotational vector fields. \pts{2}
  \item State Gauss' divergence theorem and Green's theorem in a plane. \pts{2}
  \item Define an orthogonal curvilinear coordinate system in $\mathbb{R}^3$. \pts{2}
  \item Write down the coordinate surfaces and coordinate curves for the spherical coordinate system in $\mathbb{R}^3$. \pts{2}
  \item Obtain the element of arc length for the cylindrical coordinate system in $\mathbb{R}^3$. \pts{2}
  \item Give the definition of contravariant and covariant vectors. \pts{2}
\end{parts}

\medskip\rule{\linewidth}{0.2pt}

\medskip


\noindent   \textbf{Question 2.}

\begin{parts}
  \item Prove that $
abla  \times (
abla \phi) = 0$ for any smooth scalar field $\phi$. \pts{6}
  \item Prove that $
abla \cdot (
abla  \times \vec{V}) = 0$ for any smooth vector field $\vec{V}$. \pts{6}
\end{parts}

\medskip\rule{\linewidth}{0.2pt}

\medskip


\noindent   \textbf{Question 3.}

\begin{parts}
  \item If $\vec{V} = 2y^3\vec{i} + yz^2\vec{j} + zx^2\vec{k}$, evaluate $\int_C \vec{V} \cdot d\vec{r}$ along the path $x=t$, $y=t^2$, $z=t^3$ from $t=0$ to $t=1$. \pts{6}
  \item If $\int_C \vec{V} \cdot d\vec{r}$ is independent of the path $C$ joining any two points, show that there exists a scalar function $\phi$ such that $\vec{V} = 
abla\phi$. \pts{6}
\end{parts}

\medskip\rule{\linewidth}{0.2pt}

\medskip


\noindent   \textbf{Question 4.}

\begin{parts}
  \item If $\vec{V} = y\vec{i} + (x - 2xz)\vec{j} - xy\vec{k}$, evaluate $\iint_S (
abla  \times \vec{V}) \cdot \vec{n} \, dS$, where $S$ is the surface of the sphere $x^2 + y^2 + z^2 = 1$ above the $xy$-plane. \pts{6}
  \item Using Green's theorem, show that the area bounded by a simple closed curve $C$ is $\frac{1}{2} \oint_C (x\,dy - y\,dx)$. Hence obtain the area enclosed by the ellipse $x = 2\cos \theta$, $y = 3\sin \theta$. \pts{3+3}
\end{parts}

\medskip\rule{\linewidth}{0.2pt}

\medskip


\noindent   \textbf{Question 5.}

\begin{parts}
  \item State and prove Gauss' divergence theorem and show that for a closed surface $S$:
  \[
  \iint_S \vec{r} \cdot \vec{n} \, dS = 3V
  \]
  where $V$ is the volume enclosed by $S$. \pts{6}
  \item Verify Stokes' theorem for $\vec{V} = (2x - y)\vec{i} - yz^2\vec{j} - y^2 z\vec{k}$, where $S$ is the upper hemisphere of the sphere $x^2 + y^2 + z^2 = 1$ and $C$ is its boundary. \pts{6}
\end{parts}

\medskip\rule{\linewidth}{0.2pt}

\medskip


\noindent   \textbf{Question 6.}

\begin{parts}
  \item Evaluate the triple integral:
  \[
  \iiint_V (x^2 + y^2 + z^2) \, dx\,dy\,dz
  \]
  where $V$ is the region enclosed by the sphere $x^2 + y^2 + z^2 = a^2$. \pts{6}
  \item Express $
abla  \times \vec{V}$ and $
abla^2\phi$ in spherical polar coordinates in $\mathbb{R}^3$. \pts{3+3}
\end{parts}

\medskip\rule{\linewidth}{0.2pt}

\medskip


\noindent   \textbf{Question 7.}

\begin{parts}
  \item Give the definition of Christoffel symbols of the second kind and obtain the transformation laws for them. \pts{2+4}
  \item Prove that the covariant derivatives of $g_{ij}$, $g^{ij}$, and $\delta^i_j$ are zero. \pts{2+2+2}
\end{parts}

\vfill
\rule{\linewidth}{0.4pt}

\begin{center}\small
  \href{https://bhu-exam.vercel.app/}{Click here to visit our website}\\[4pt]
  \href{https://me-aryan.vercel.app/}{Click here to visit our contributor}\\[4pt]
  \href{https://quashan-me.vercel.app/}{Click here to visit our contributor}
\end{center}

\end{document}
